\scanpage{15}
\begin{center}\small Sur les groupes de Galois attachés aux groupes $p$-divisibles\end{center}
\mspage{129}
Soit $\bar A$ (resp. $\overline{\mathfrak{m}}$) l’anneau (resp. l’idéal) de valuation de $\bar K$. Pour tout entier $n\geqslant 0$, notons $f^{(n)}$ le $n$-ème itéré de $f$ (i.e. \textit{la multiplication par $p^n$} dans $F$), et soit $T_n$ l’ensemble des $x\in\overline{\mathfrak{m}}$ tels que $f^{(n)}(x)=0$. On sait (cf. par exemple [\textit{3}]) que $T_n$ est un groupe (pour la loi de groupe formel définissant $F$), et que ce groupe est isomorphe à $(\mathbf{Z}/p^n\mathbf{Z})^h$; on a $T=\varprojlim T_n$. Notons $T'_n$ l’ensemble des éléments de $T_n$ d’ordre $p^n$; on a
\[
T'_n=T_n-T_{n-1}.
\]

\textbf{Lemme 3.} -- \textit{Il existe un nombre $c>0$ tel que, pour tout $n\geqslant 1$, et tout $x\in T'_n$, l’indice de ramification de l’extension $K(x)/K$ soit $\geqslant c\cdot p^{nh}$.}

Admettons provisoirement ce lemme, et démontrons la prop.~8. Quitte à transformer $z$ par une homothétie, on peut supposer que $z\in T$ et $z\notin pT$. Pour tout $n\geqslant 1$, l’image $z_n$ de $z$ dans $T/p^nT=T_n$ appartient à $T'_n$. Soit d’autre part $U=G\cdot z$ \textit{l’orbite} de $z$ par $G$. C’est une \textit{sous-variété analytique $p$-adique} de $T$ et son espace tangent au point $z$ est égal à $\mathfrak{g}\cdot z$ (cf. par exemple [\textit{7}], LG, p.~4.12). D’autre part, l’image $U_n$ de $U$ dans $T/p^nT$ est égale à $G\cdot z_n$, ensemble des conjugués de $z_n$. Mais, d’après le lemme 3, on a $\operatorname{Card}(G\cdot z_n)=[K(z_n):K]\geqslant c\cdot p^{nh}$. D’où
\[
\operatorname{Card}(U_n)\geqslant c\cdot p^{nh}.
\]
Soit alors $\mu$ la mesure de Haar de $T$, normalisée de telle sorte que $\mu(T)=1$. La formule ci-dessus entraîne que $\mu(U)\geqslant c$. Mais il est immédiat qu’une sous-variété analytique $p$-adique de $T$ n’a une mesure $>0$ que si son intérieur est non vide. Comme $U$ est une orbite, on en conclut que $U$ est \textit{ouvert} dans $T$. Son espace tangent en $z$ est donc $V$ tout entier. D’où $\mathfrak{g}\cdot z=V$.

\textit{Reste à démontrer le lemme} 3. Soit $v$ la valuation de $\bar K$, normalisée de telle sorte que $v(p)=1$. Soit $c_1$ la borne inférieure des $v(a_n)$, pour $n<p^h$. D’après ce qui a été dit plus haut, on a $c_1>0$. Soit $\varphi:\mathbf{R}_+\to\mathbf{R}_+$ l’application donnée par la formule
\[
\varphi(\alpha)=\operatorname{Inf}(p^h\alpha,\alpha+c_1).
\]
C’est une bijection strictement croissante de $\mathbf{R}_+$ sur $\mathbf{R}_+$; soit $\psi$ la bijection réciproque. Si $x\in\overline{\mathfrak{m}}$, la formule
\[
f(x)=\sum a_nx^n
\]
montre que
\[
v(f(x))\geqslant\varphi(v(x)),
\]
% archive sheet number: 14
